\subsection{Arquitectura del Módulo de Autenticación}
 
El módulo de autenticación es el punto de entrada transversal de la
plataforma: cualquier operación protegida —registro de intentos, consulta
de alumnos, asignación de ejercicios— depende de que el usuario haya
obtenido previamente un token válido. El diseño sigue el patrón
arquitectónico MVC establecido para el sistema y se apoya en el estándar
JSON Web Token (JWT) para la transmisión segura de identidad y rol entre
el Front-End y el Back-End.
 
La arquitectura del módulo se divide en tres componentes principales:
 
\begin{enumerate}
    \item Capa de Presentación (Front-End): Compuesta por los
    componentes \texttt{Login.jsx} y \texttt{Registro.jsx}. Esta capa es
    responsable de la validación de formato en el cliente (longitud de
    campos, caracteres permitidos), el envío de las credenciales al
    Back-End mediante peticiones HTTP y el almacenamiento del token
    recibido en \texttt{localStorage} para su uso en peticiones
    subsecuentes.
 
    \item Capa de Lógica de Negocio (Back-End): Implementada
    en el controlador \texttt{auth.py} mediante FastAPI. Esta capa
    gestiona la verificación de credenciales contra la base de datos,
    el cifrado y la validación de contraseñas con bcrypt, la emisión de
    tokens JWT firmados y la verificación del estado activo del usuario.
    Adicionalmente, el módulo \texttt{dependencies.py} expone las
    dependencias reutilizables \texttt{require\_tutor} y
    \texttt{require\_alumno}, que actúan como guardas de acceso en todos
    los demás controladores del sistema.
 
    \item Capa de Persistencia (Base de Datos): El módulo opera
    exclusivamente sobre la tabla \texttt{Usuario}, consultando el campo
    \texttt{contrasena\_cifrada} (almacenado como hash bcrypt en
    \texttt{VARCHAR(255)}) y el campo \texttt{id\_estatus} para verificar
    que la cuenta se encuentre activa antes de emitir el token.
\end{enumerate}
 
% ─────────────────────────────────────────────────────────
\subsubsection{Diseño del Token JWT}
 
El token emitido por el sistema sigue el estándar JWT (RFC 7519) firmado
con el algoritmo simétrico HMAC-SHA256 (HS256). El payload del token
contiene exclusivamente los campos necesarios para que cualquier
controlador pueda tomar decisiones de autorización sin consultar la base
de datos en cada petición:
 
\begin{lstlisting}[language=json, caption={Estructura del payload JWT emitido por el Back-End}, label={lst:jwt-payload}]
{
  "id_usuario": <integer>,
  "nombre":     <string>,
  "tipo_usuario": "tutor" | "alumno",
  "exp":        <unix timestamp>
}
\end{lstlisting}
 
El campo \texttt{tipo\_usuario} es el discriminador de rol que consumen
las dependencias \texttt{require\_tutor} y \texttt{require\_alumno}.
El token tiene una vigencia de 24 horas (\texttt{ACCESS\_TOKEN\_EXPIRE\_MINUTES = 1440}),
tras la cual el usuario debe autenticarse nuevamente. La clave de firma
(\texttt{SECRET\_KEY}) se gestiona mediante variables de entorno y no se
incluye en el repositorio.
 
% ─────────────────────────────────────────────────────────
\subsubsection{Diseño del Flujo de Interacción}
 
El módulo contempla dos flujos principales: inicio de sesión y registro
de tutor.
 
\paragraph{Flujo de Inicio de Sesión:}
\begin{enumerate}
    \item El usuario introduce su \texttt{username} y contraseña en el
    componente \texttt{Login.jsx} y envía la petición.
    \item El Front-End realiza una petición \texttt{POST} al endpoint
    \texttt{/api/auth/login} con el payload en formato JSON.
    \item El Back-End consulta la tabla \texttt{Usuario} filtrando por
    \texttt{username}. Si el registro no existe, retorna un error
    \texttt{401 Unauthorized} con un mensaje genérico para evitar la
    enumeración de usuarios.
    \item Se verifica que \texttt{id\_estatus = 1} (cuenta activa). Si
    el valor es distinto, se retorna \texttt{403 Forbidden}.
    \item Se verifica la contraseña ingresada contra el hash almacenado
    utilizando bcrypt. Un fallo retorna \texttt{401 Unauthorized} con
    el mismo mensaje genérico del paso 3.
    \item Si todas las validaciones son exitosas, se genera el token JWT
    con el payload descrito en la sección anterior y se retorna al
    Front-End junto con un mensaje de confirmación.
    \item El Front-End almacena el token en \texttt{localStorage} y
    redirige al usuario a su panel correspondiente según el campo
    \texttt{tipo\_usuario} incluido en el payload decodificado.
\end{enumerate}
 
\paragraph{Flujo de Registro de Tutor:}
\begin{enumerate}
    \item El tutor introduce sus datos en el componente
    \texttt{Registro.jsx}. El Front-End aplica validaciones de formato
    antes del envío.
    \item El Front-End realiza una petición \texttt{POST} al endpoint
    \texttt{/api/auth/register} con el payload en formato JSON.
    \item El Back-End verifica que el correo electrónico y el
    \texttt{username} no existan previamente en la tabla
    \texttt{Usuario}. En caso de duplicado, retorna \texttt{400 Bad Request}
    con el campo en conflicto especificado.
    \item La contraseña es validada contra las reglas de complejidad
    definidas en el modelo Pydantic (longitud mínima de 8 caracteres,
    al menos una mayúscula, una minúscula, un dígito y un carácter
    especial). Una contraseña inválida retorna \texttt{422 Unprocessable
    Entity}.
    \item La contraseña válida se cifra con bcrypt y se inserta un nuevo
    registro en la tabla \texttt{Usuario} con \texttt{tipo\_usuario =
    'tutor'} e \texttt{id\_estatus = 1} (activo por defecto).
    \item El Back-End retorna \texttt{200 OK} con un mensaje de
    confirmación. El tutor debe iniciar sesión de forma explícita
    para obtener su token.
\end{enumerate}
 
\noindent \textit{Nota de diseño:} el registro está restringido a cuentas
de tipo tutor. Los alumnos son dados de alta únicamente por su tutor
asignado a través del módulo de gestión de alumnos, lo que garantiza que
ningún alumno pueda auto-registrarse en la plataforma.
 
% ─────────────────────────────────────────────────────────
\subsubsection{Diseño de APIs}
 
Para soportar los flujos descritos, el diseño de la API contempla los
siguientes endpoints en el controlador \texttt{auth.py}:
 
\textbf{Endpoint 1:} Inicio de Sesión
 
El endpoint \texttt{/api/auth/login} opera mediante el método
\texttt{POST} y no requiere autenticación previa. Recibe un payload
JSON con \texttt{username} y \texttt{password} y, ante credenciales
válidas, retorna el token JWT.
 
\begin{lstlisting}[language=json, caption={Payload de solicitud y respuesta del endpoint de login}, label={lst:login-payload}]
/* Solicitud */
{
  "username": "string",
  "password": "string"
}
 
/* Respuesta 200 OK */
{
  "message": "Inicio de sesion exitoso",
  "token":   "eyJhbGci..."
}
\end{lstlisting}
 
\textbf{Endpoint 2:} Registro de Tutor
 
El endpoint \texttt{/api/auth/register} opera mediante el método
\texttt{POST} y no requiere autenticación previa. Recibe los datos del
nuevo tutor y aplica las validaciones de unicidad y complejidad de
contraseña descritas en el flujo anterior.
 
\begin{lstlisting}[language=json, caption={Payload de solicitud del endpoint de registro}, label={lst:register-payload}]
/* Solicitud */
{
  "nombre":   "string",
  "username": "string (min 6 chars, [a-zA-Z0-9_])",
  "email":    "string",
  "password": "string (min 8 chars, mayus + minus + num + especial)"
}
 
/* Respuesta 200 OK */
{
  "message": "Usuario registrado exitosamente"
}
\end{lstlisting}
 
% ─────────────────────────────────────────────────────────
\subsubsection{Diseño del Control de Acceso Basado en Roles}
 
El sistema implementa un control de acceso basado en roles (RBAC) a
través de dos dependencias reutilizables definidas en
\texttt{dependencies.py}: \texttt{require\_tutor} y
\texttt{require\_alumno}. Ambas se inyectan como parámetros en los
endpoints que requieren autenticación, siguiendo el mecanismo de
inyección de dependencias de FastAPI.
 
El flujo de validación es el siguiente:
 
\begin{enumerate}
    \item El cliente incluye el token en el encabezado HTTP de cada
    petición protegida: \texttt{Authorization: Bearer <token>}.
    \item La dependencia \texttt{get\_current\_user} extrae y decodifica
    el token usando la misma \texttt{SECRET\_KEY} y algoritmo HS256. Si
    el token está vencido o es inválido, retorna \texttt{401 Unauthorized}.
    \item Las dependencias \texttt{require\_tutor} y
    \texttt{require\_alumno} verifican el campo \texttt{tipo\_usuario}
    del payload decodificado. Si el rol no corresponde al requerido por
    el endpoint, retornan \texttt{403 Forbidden}.
    \item Si la validación es exitosa, el objeto \texttt{current\_user}
    (con \texttt{id\_usuario}, \texttt{tipo\_usuario} y \texttt{nombre})
    queda disponible para la lógica del controlador sin necesidad de
    realizar consultas adicionales a la base de datos.
\end{enumerate}
 
La Tabla~\ref{tab:endpoints-auth} resume los endpoints del sistema
clasificados por el tipo de autenticación que requieren.
 
\begin{table}[H]
\centering
\caption{Clasificación de endpoints por nivel de acceso}
\label{tab:endpoints-auth}
\begin{tabular}{|l|l|l|}
\hline
\textbf{Endpoint} & \textbf{Método} & \textbf{Acceso requerido} \\
\hline
\texttt{/api/auth/login}              & POST & Público \\
\texttt{/api/auth/register}           & POST & Público \\
\texttt{/api/intentos/registrar}      & POST & Alumno (JWT) \\
\texttt{/api/intentos/tutor/\{id\}}   & GET  & Tutor (JWT) \\
\texttt{/api/intentos/calificar/\{id\}} & PUT & Tutor (JWT) \\
\texttt{/api/alumnos/...}             & *    & Tutor (JWT) \\
\texttt{/api/ejercicios/...}          & *    & Tutor / Alumno (JWT) \\
\texttt{/api/perfil/...}              & *    & Tutor / Alumno (JWT) \\
\hline
\end{tabular}
\end{table}